%% ma: L12-B15-0016
%% lop: 12
%% chuong: 5
%% bai: 15
%% dang: 12.15.13
%% loai: DS
%% muc: [NB, TH, TH, VD]
%% dap_an: "SSSS"
%% nguon: "(NBV)-10 ĐỀ KIỂM TRA HK2-TOÁN 12-DUNG CHUNG 3 SÁCH.pdf (H01, Đề số 1), Phần 2 Câu 2"
%% kien_thuc: "Bài 15 (vectơ chỉ phương của đường thẳng), Bài 14 (vectơ pháp tuyến, tích có hướng), đường thẳng vuông góc với mặt phẳng, hình chiếu lên trục"
%% trang_thai: chinh-thuc
%% ghi_chu: "Đáp án gốc ghi Đ Đ S S nhưng chính lời giải gốc kết luận a) vectơ chỉ phương của Δ1 là (0;-3;4) (đề in (1;-3;4), nên a) sai) và b) không vuông góc (hai vectơ không cùng phương, nên b) sai). Dùng S S S S (giải độc lập: c) [v;k] = (-3;-3;0) ∥ (1;1;0), không phải (3;-3;2); d) H(0;0;2), AH = (-1;1;0), không phải (-1;-1;0)). Nếu đề gốc muốn ý a) là (0;-3;4) thì a) đúng. Đã duyệt 19/10/2026."
\begin{ex}
Trong không gian $Oxyz$, cho các đường thẳng $\Delta_1\colon\begin{cases}x=1\\y=2-3t\\z=3+4t\end{cases}$, $\Delta_2\colon\dfrac{x-1}3=\dfrac y{-3}=\dfrac{z+3}2$ và mặt phẳng $(P)\colon x+3y-2z+1=0$.
\choiceTF
{Vectơ chỉ phương của đường thẳng $\Delta_1$ là $\vec a=(1;-3;4)$.}
{Đường thẳng $\Delta_1$ vuông góc với $(P)$.}
{Đường thẳng $d_2$ vuông góc với $\Delta_2$ và song song với $(Oxy)$ thì có một vectơ chỉ phương là $\vec u_2=(3;-3;2)$.}
{Đường thẳng $d_3$ qua $A(1;-1;2)$, cắt và vuông góc với trục $Oz$ thì có vectơ chỉ phương là $\vec u_3=(-1;-1;0)$.}
\loigiai{a) Từ phương trình tham số, $\Delta_1$ có vectơ chỉ phương $(0;-3;4)$, không phải $(1;-3;4)$: a) sai.

b) Vectơ pháp tuyến của $(P)$ là $\vec n=(1;3;-2)$. Vì $(0;-3;4)$ và $(1;3;-2)$ không cùng phương nên $\Delta_1$ không vuông góc với $(P)$: b) sai.

c) $\Delta_2$ có vectơ chỉ phương $\vec v=(3;-3;2)$, $(Oxy)$ có vectơ pháp tuyến $\vec k=(0;0;1)$. Đường thẳng $d_2$ có vectơ chỉ phương $[\vec v,\vec k]=(-3;-3;0)$, cùng phương với $(1;1;0)$, không phải $(3;-3;2)$ (vectơ này không vuông góc với $\vec v$): c) sai.

d) Gọi $H=d_3\cap Oz$. Vì $d_3\perp Oz$ nên $H$ là hình chiếu vuông góc của $A$ lên $Oz$, $H(0;0;2)$. Do đó $d_3$ có vectơ chỉ phương $\overrightarrow{AH}=(-1;1;0)$, không cùng phương với $(-1;-1;0)$: d) sai.}
\end{ex}
